\documentclass{article}
\usepackage{fontspec}
\usepackage[UTF8, scheme=plain]{ctex}
\usepackage{emoji}
\usepackage{geometry}
\geometry{a4paper, margin=1in}

\title{CS336 Assignment 1 Basics Writeup}
\author{}
\date{}

\begin{document}

\maketitle

\section*{Problem (unicode1): Understanding Unicode}

\begin{enumerate}
    \item[(a)] \texttt{chr(0)} returns the Null character (U+0000).
    \item[(b)] The string representation (\texttt{\_\_repr\_\_()}) displays the escaped hex code \texttt{'\\x00'}, whereas the printed representation outputs the actual invisible control character.
    \item[(c)] When this character occurs in a Python string, it behaves as a normal character (increasing length without terminating the string) but is typically invisible when printed to the console.
\end{enumerate}

\section*{Problem (unicode2): Unicode Encodings}

\begin{enumerate}
    \item[(a)] UTF-8 is preferred because it is variable-width and highly space-efficient for ASCII characters (using 1 byte vs 2-4 bytes in UTF-16/32), which dominate most training corpora (like English text and code), and it avoids the issue of null-byte sparsity found in fixed-width encodings.
    \item[(b)] "A中\emoji{grinning-face}". The function fails to decode the bytes correctly because it does not handle multi-byte characters properly. It attempts to decode each byte as a separate character, which can lead to errors or unexpected results.
    \item[(c)] \texttt{"\textbackslash xe4\textbackslash x00"}. The sequence starts with \texttt{"\textbackslash xe4"}(11100100), implicating that it is a 3-byte sequence. However, the second byte is \texttt{"\textbackslash x00"}(00000000) instead of "10xxxxxx", so it is not a valid UTF-8 encoding of a character, and there are only two bytes available to decode.
\end{enumerate}

\section*{Problem (train\_bpe): BPE Tokenizer Training}

\begin{enumerate}
    \item[] Implemented \texttt{pretokenize}, \texttt{init\_vocab}, \texttt{build\_words}, \texttt{init\_pairs}, \texttt{merge}.
\end{enumerate}

\section*{Problem (train\_bpe\_tinystories): BPE Training on TinyStories}

\begin{enumerate}
    \item[(a)] 320s in 4vcpu, 38.09s in 128vcpu, both took up about 3.4 GB RAM. Longest token (ID 7160): b' accomplishment'. 
    \item[(b)] Pretokenize took the most time.
\end{enumerate}

\section*{Problem (train\_bpe\_expts\_owt): BPE Training on OpenWebText}

\begin{enumerate}
    \item[(a)] 3682.80s in 128vcpu. Longest token (ID 25851): b'\xc3\x83\xc3\x82\xc3\x83\xc3\x82\xc3\x83\xc3\x82\xc3\x83\xc3\x82\xc3\x83\xc3\x82\xc3\x83\xc3\x82\xc3\x83\xc3\x82\xc3\x83\xc3\x82\xc3\x83\xc3\x82\xc3\x83\xc3\x82\xc3\x83\xc3\x82\xc3\x83\xc3\x82\xc3\x83\xc3\x82\xc3\x83\xc3\x82\xc3\x83\xc3\x82\xc3\x83\xc3\x82'. 
    \item[(b)] Pretokenize took the most time.
\end{enumerate}

\end{document}
